\documentclass[10pt,a4paper]{amsart}
\usepackage[margin=19mm]{geometry}
\usepackage{amsmath,amssymb,mathtools}
\usepackage[scr=rsfs,cal=euler]{mathalfa}
\usepackage{xfrac}
\usepackage[unicode,hidelinks,bookmarks=false]{hyperref}
\newcommand{\ie}{i.e.\ }
\newcommand{\cf}{cf.\ }
\newcommand{\resp}{respectively\ }
\newcommand{\Subsection}{\subsection}
\providecommand{\nobreakdash}{\nobreak\hbox{-}}
\setlength{\emergencystretch}{2em}
\multlinegap0pt
\usepackage{bm}
\usepackage{bbm}
\usepackage{dsfont}
\usepackage{xspace}
\usepackage{cancel}
\usepackage{mathtools}
\usepackage{amsthm}
\makeatletter
\@ifundefined{thm}{\newtheorem{thm}{Thm}}{}
\@ifundefined{lem}{\newtheorem{lem}[thm]{Lemma}}{}
\makeatother
\makeatletter
\expandafter\ifx\csname enumerateA\endcsname\relax
\newenvironment{enumerateA}{\begin{enumerate}[\upshape (A)]}{\end{enumerate}}
\fi
\expandafter\ifx\csname enumeratea\endcsname\relax
\newenvironment{enumeratea}{\begin{enumerate}[\upshape (a)]}{\end{enumerate}}
\fi
\expandafter\ifx\csname enumeratei\endcsname\relax
\newenvironment{enumeratei}{\begin{enumerate}[\upshape (i)]}{\end{enumerate}}
\fi
\expandafter\ifx\csname enumeraten\endcsname\relax
\newenvironment{enumeraten}{\begin{enumerate}[\upshape 1.]}{\end{enumerate}}
\fi
\renewcommand{\geq}{\geqslant} 
\let\gt=\tau \let\gth=\vartheta
\newcommand{\la}{\langle} 
\renewcommand{\le}{\leqslant} 
\renewcommand{\leq}{\leqslant} 
\newcommand{\mvgs}{{\sigma}}
\newcommand{\mvgt}{{\tau}}
\newcommand{\ra}{\rangle}
\makeatother
\title[Joint parameter estimations for spin glasses]{Joint parameter estimations for spin glasses}
\author{Wei-Kuo Chen, Arnab Sen, Qiang Wu}
\date{Source: arXiv:2406.10760v2 (CC BY 4.0)}
\begin{document}
\maketitle

\noindent\emph{Source and licence.} Joint parameter estimations for spin glasses. Version arXiv:2406.10760v2 (CC BY 4.0), distributed under the Creative Commons Attribution 4.0 International licence, as recorded in the arXiv licence metadata for this exact version and shown on the abstract page https://arxiv.org/abs/2406.10760v2. Full rights notice: licenses/arxiv-2406.10760-CC-BY-4.0.txt.\par\smallskip

\noindent\textit{Editorial scope.} Counted unit: the Lem whose source label is \texttt{lem:new-concent} in the article (source0.tex lines 884--892, source label \texttt{lem:new-concent}), delivered together with the article's own complete original proof of it (source0.tex lines 895--998, one \texttt{proof} environment of 104 lines). No mathematics is written, repaired or re-derived: statement and proof are the article's own text line for line. Required context retained uncounted: none. Every definition, display and label of those lines is retained unchanged. Omitted from this package and not redistributed: 1--883, 893--894, 999--2288. Nothing inside a retained excerpt is omitted or altered: every retained body is delivered exactly as the article's lines are written, so no omission receipt of any kind accompanies this package. Citations: the retained excerpts carry no citation command at all, so no bibliography entry of the article is transcribed and no reference is redistributed. Reference adaptation: none was needed and none was made; every reference inside the delivered unit resolves inside it, so no unresolved reference or citation is produced. Printed slips kept literal: no printed word, symbol, hypothesis or number of the counted statement or of its proof was altered. Layout and class: the article's own class is \texttt{amsart} with packages including geometry, systeme, mathrsfs, xfrac, hyperref, amsmath, amssymb, amsthm, amsfonts, amsbsy, latexsym, dsfont, color, amsrefs; the wrapper is the lane's pinned \texttt{amsart} editorial wrapper at 10pt on A4, which restates the theorem-like environments the retained text needs and adds the article's own macro definitions that the retained text needs (\texttt{\textbackslash geq}, \texttt{\textbackslash gt}, \texttt{\textbackslash la}, \texttt{\textbackslash le}, \texttt{\textbackslash leq}, \texttt{\textbackslash mvgs}, \texttt{\textbackslash mvgt}, \texttt{\textbackslash ra}), with extra wrapper packages (bm, bbm, dsfont, xspace, cancel, mathtools). The delivered numbering is the unit's own. Assets: no retained span contains a figure environment, a graphics-inclusion command, a PDF-page inclusion, a table or a TikZ picture; no image file is copied into this package, the package declares no asset dependency and no image file is redistributed with it. Licence: version arXiv:2406.10760v2 of this article is distributed under the Creative Commons Attribution 4.0 International licence, as recorded in the arXiv licence metadata for this exact version and shown on the abstract page https://arxiv.org/abs/2406.10760v2. A full rights notice is delivered at licenses/arxiv-2406.10760-CC-BY-4.0.txt. Characters: every retained excerpt is pure ASCII, so the delivered engine, XeLaTeX with its default Latin Modern text font, reports no missing character.
\begin{lem} \label{lem:new-concent}
Let $\mvgt$ be sampled from $P_{\beta, h}$ for some $\beta>0$ and $h\in \mathbb{R}$.
Let $\{f_i\}_{i\in [N]}$ be a collection of functions defined on $\{-1,+1\}\times \mathbb{R}$ with $\max_{i\in [N]}|f_i(x,y)|\leq C(1+|y|)$ for all $(x,y)\in \{-1,+1\}\times\mathbb{R}$ for some positive constant $C$ and $\max_{i\in [N]}(\|\partial_yf_i\|_\infty,\|\partial_{y}^2f_i\|_\infty)<\infty.$
Then there exists a constant $K$ depending only on $\beta$ and $C$ such that for any $N\geq 1,$ we have that 
    \begin{align*}
&\Bigl\la\Bigl(\sum_{i=1}^N\bigl(f_i(\tau_i,m_i(\mvgt))-L_i(\mvgt)\bigr)\Bigr)^2 \Bigr\ra\\
&\leq K(1+\|J\|+\|J\|^2)\Bigl(N+N\|J\|^2+\sum_{i=1}^N\|\partial_y f_i\|_\infty^2+\sum_{i=1}^N\|\partial_y^2f_i\|_\infty^2\Bigr).
\end{align*}
\end{lem}

\begin{proof}
For any $i,j\in[N]$, set 
$$
m_j^{(i)}(\mvgs):=\sum_{k\neq i, j}J_{jk}\sigma_k
$$ 
and let $L^{(i)}_j(\mvgs)$ be the same as $L_j(\mvgs)$ with $m_j(\mvgs)$ replaced by $m_j^{(i)}(\mvgs)$.
Since $L_i(\tau)$ is the expectation of $f_i(\tau_i, m_i(\mvgt))$ conditionally on $(\tau_l)_{l\neq i}$, we have
\begin{align*}
    \bigl\la\bigl(f_i(\tau_i,m_i(\mvgt))-L_i(\mvgt)\bigr)
    \bigl(f_j(\tau_j,m_j^{(i)}(\mvgt))-L_j^{(i)}(\mvgt)\bigr)\bigr\ra=0,
\end{align*}
which implies
\begin{align}
   &  \bigl\la \bigl(f_i(\tau_i, m_i(\mvgt) )-L_i(\mvgt)\bigr)\bigl(f_j(\tau_j, m_j(\mvgt) )-L_j(\mvgt)\bigr)\bigr\ra  \nonumber\\
    & = \bigl\la  X_i(\tau) \big (f_j(\tau_j, m_j(\mvgt))-f_j(\tau_j, m_j^{(i)}(\mvgt))\big ) \bigr\ra \label{eq:ip_decom1}\\
    & + \bigl\la X_i(\tau)c_{ij}(\mvgt)\bigr\ra, \label{eq:ip_decom2}
\end{align}
where
\begin{align*}
c_{ij}(\mvgt)&:=L_j^{(i)}(\mvgt)-L_j(\mvgt),\\
X_i(\gt)&:= f_i(\tau_i,  m_i(\mvgt) )-L_i(\mvgt).
\end{align*}
We bound \eqref{eq:ip_decom2} first.
Note that 
\begin{align}\label{ProofofCor2.2:eq0}
|X_i(\gt)| \le 2 C(1+|m_i(\tau)|).
\end{align}
For $t\in [0,1],$ let $J^{(i,j)}(t)$ be a symmetric square matrix with $J_{ab}^{(i,j)}(t)=J_{ab}$ for all $(a,b)\neq (i,j)$ and $(j,i)$ and $J_{ij}^{(i,j)}(t)=J_{ji}^{(i,j)}(t)=(1-t)J_{ij}.$ Denote $\mathcal{L}_{ij}(t)$ by the function $L_j(\tau)$ with $J$ being replaced by $J^{(i,j)}(t)$. Using Taylor's formula yields that
\begin{align*}
    L_j^{(i)}(\mvgt)&=\mathcal{L}_{ij}(1)\\
    &=\mathcal{L}_{ij}(0)+\mathcal{L}_{ij}'(0)+\frac{1}{2}\int_0^1(1-s)\mathcal{L}_{ij}''(s)ds\\
    &=L_j(\tau)+J_{ij}Y_j+J_{ij}^2Z_j^{(i)},
\end{align*}
where
\begin{align*}
    Y_j&=-\partial_{J_{ij}}L_j(\tau),\\
    Z_j^{(i)}&=\frac{1}{2}\int_0^1(1-s)\partial_{J_{ij}}^2L_j(\tau)\Big|_{J=J^{(i,j)}(s)}ds.
\end{align*}
We emphasize that here $Y_j$ depends only on $j$ since
\begin{align*}
    \partial_{J_{ij}}L_j(\mvgt)&=\frac{e^{\beta m_j(\mvgt)+h}\partial_yf_j(1,m_j(\mvgt))+e^{-\beta m_j(\mvgs)-h}\partial_yf_j(-1,m_j(\mvgt))}{e^{\beta m_j(\mvgt)+h}+e^{-\beta m_j(\mvgt)-h}}\\
    &+\beta\frac{e^{\beta m_j(\mvgt)+h}f_j(1,m_j(\mvgt))-e^{-\beta m_j(\mvgt)-h}f_j(-1,m_j(\mvgt))}{e^{\beta m_j(\mvgt)+h}+e^{-\beta m_j(\mvgt)-h}}\\
    &-\beta\frac{e^{\beta m_j(\mvgt)+h}f_j(1,m_j(\mvgt))+e^{-\beta m_j(\mvgt)-h}f_j(-1,m_j(\mvgt))}{e^{\beta m_j(\mvgt)+h}+e^{-\beta m_j(\mvgt)-h}}\tanh(\beta m_j(\mvgt)+h).
\end{align*}
Now from the given assumption,
\begin{align*}
  |\partial_{J_{ij}}L_j(\mvgt)|&\leq 2\beta C(1+|m_j(\tau)|)+\|\partial_yf_j\|_\infty=:K_j.
\end{align*}
Similarly one can compute the second derivative of $L_j(\mvgt)$ with respect to $J_{ij}$; the exact form is not important here, but it is crucial to realize that
\begin{align*}
    |\partial_{J_{ij}}^2L_j(\mvgt)|&\leq \tilde C\bigl(1+|m_j(\tau)|+\|\partial_yf_j\|_\infty+\|\partial_y^2f_j\|_\infty\bigr)
\end{align*}
for some constant $\tilde C$ depending only on $C$ and $\beta.$ It follows from these bounds that 
\begin{align}\label{ProofofCor2.2:eq1}
    |Y_j|\leq K_j
\end{align} and
\begin{align}
 \nonumber   |Z_j^{(i)}|&\leq \tilde{C}\Bigl(1+\int_0^1(1-s)|(J^{(i,j)}(s)\tau)_j\bigr|ds+\|\partial_yf_j\|_\infty+\|\partial_y^2f_j\|_\infty\Bigr)\\
\label{ProofofCor2.2:eq2}    &\leq \tilde{C}\bigl(1+2\|J\|+\|\partial_yf_j\|_\infty+\|\partial_y^2f_j\|_\infty\bigr)=:\tilde{K}_j,
\end{align}
where we have used the bound $\|J^{(i,j)}(s)\|\leq \|J\|+(1-s)|J_{ij}|\leq 2\|J\|.$
Now, write
\[ \Big| \sum_{i \ne j} X_i(\tau) c_{ij}(\tau) \Big |\le \Big| \sum_{i \ne j} J_{ij} X_i(\tau) Y_j(\tau) \Big| + \sum_{i \ne j} \frac{J_{ij}^2}{2}|X_i(\tau)| |Z_j^{(i)}| .
\]
Note that $J_{ij}=0$ when $i=j$. Using \eqref{ProofofCor2.2:eq0} and \eqref{ProofofCor2.2:eq1} leads to
\begin{align*}
\Bigl|  \sum_{i\neq j} J_{ij} X_i(\gt) Y_j(\tau)  \Bigr|&\leq 2C\|J\|\Bigl(\sum_{i=1}^N(1+|m_i(\tau)|)^2\Bigr)^{1/2}\Bigl(\sum_{j=1}^NK_j^2\Bigr)^{1/2}.
\end{align*}
Similarly, from \eqref{ProofofCor2.2:eq0} and \eqref{ProofofCor2.2:eq2} and using $\|J\circ J\|\leq \|J\|^2,$
\begin{align*}
     \sum_{i \ne j} J_{ij}^2|X_i(\tau)||Z_j^{(i)}|  
    &\leq C\|J\|^2\Bigl(\sum_{i=1}^N(1+|m_i(\tau)|)^2\Bigr)^{1/2}\Bigl(\sum_{i=1}^N\tilde K_i^2\Bigr)^{1/2}.
\end{align*}
From these, \eqref{eq:ip_decom2} is bounded by
\begin{align*}
  \Big| \sum_{i \ne j} X_i(\tau) c_{ij}(\tau) \Big |&\leq 2 C\|J\|\Bigl(\sum_{i=1}^N(1+|m_i(\tau)|)^2\Bigr)^{1/2}\Bigl(\sum_{j=1}^NK_j^2\Bigr)^{1/2}\\
&+C\|J\|^2\Bigl(\sum_{i=1}^N(1+|m_i(\tau)|)^2\Bigr)^{1/2}\Bigl(\sum_{i=1}^N\tilde K_i^2\Bigr)^{1/2}.
\end{align*}


As for \eqref{eq:ip_decom1}, using the Taylor expansion yields
\begin{align*}
    \Bigl|f_i(\tau_j,m_j(\mvgt))-f_i(\tau_j,m_j^{(i)}(\mvgt))+J_{ij}\tau_i\partial_yf_i(\tau_j, m_j(\mvgt))\Bigr|&\leq \frac{1}{2}J_{ij}^2\|\partial^2_{y}f_i\|_\infty.
\end{align*}
%\begin{align*}
%    \Bigl|f(\tau_j,\tau_j m_j^{(i)}(\mvgt)) - f(\tau_j,\tau_j m_j(\mvgt))+J_{ij}\tau_i\tau_j\partial_yf(\tau_j,\tau_j m_j(\mvgt))\Bigr|&\leq \frac{1}{2}J_{ij}^2\|\partial^2_{y}f\|_\infty.
%\end{align*}
Here, we can control similar to the above,
\begin{align*}
   &\Bigl| \sum_{i\neq j}X_i(\tau)\bigl(f_j(\tau_j, m_j(\mvgt))-f_j(\tau_j,m_j^{(i)})\bigr)\Bigr| \\
   &\le \Big| \sum_{i, j} J_{ij} (X_i(\tau) \tau_i ) \big(\partial_yf_j(\tau_j,m_j(\mvgt)) \big) \Big|  + \frac12  \sum_{i, j} J_{ij}^2|X_i(\gt)|\|\partial^2_y f_j\|_{\infty}  \\
   &\leq 2C\|J\|\Bigl(\sum_{i=1}^N(1+|m_i(\tau)|)^2\Bigr)^{1/2}\Bigl(\sum_{i=1}^N\|\partial_yf_i\|_\infty^2\Bigr)^{1/2}\\
&+C\|J\|^2\Bigl(\sum_{i=1}^N(1+|m_i(\tau)|)^2\Bigr)^{1/2}\Bigl(\sum_{i=1}^N\|\partial_y^2f_i\|_\infty^2\Bigr)^{1/2}.
\end{align*}
Putting everything together and noting that $\sum_{i=1}^Nm_i(\tau)^2\leq N\|J\|^2$, we have
\begin{align*}
   &\sum_{i\neq j}\bigl\la \bigl(f_i(\tau_i, m_i(\mvgt) )-L_i(\mvgt)\bigr)\bigl(f_i(\tau_j, m_j(\mvgt) )-L_j(\mvgt)\bigr)\bigr\ra\\
   &\leq K(\|J\|+\|J\|^2)\Bigl(N+N\|J\|^2+\sum_{i=1}^N\|\partial_y f_i\|_\infty^2+\sum_{i=1}^N\|\partial_y^2f_i\|_\infty^2\Bigr)
\end{align*}
for some constant $K$ depending only on $\beta$ and $C.$ 
Finally, for the sum over $i=j$, it can be easily bounded by $C''N(+\|J\|^2)$ for some constant $K'$ depending only on $C$ and this completes our proof.


\end{proof}

\end{document}
